% !PRINTED 713
\weberpage{713}
Both follow easily from the fact that for $n = -1$ in the first
case $\gamma = \tfrac12\mu$, and in the second $\alpha=1$,
$\beta=0$, so that in both cases $\chi_1(-1)=1$, and from the fact
that each of the formulae (6), (7) represents exactly
$\tfrac12\mu$ different characters.

What now matters essentially is the following:

\begin{quote}
\gesperrt{If $p$ is any prime different from $q$, then in the form
$\chi_1(p)$, as $\chi_1$ runs through the group $C_1$, one obtains
every $f_1$th root of unity, and each equally often, that is,
$e_1$ times.}
\end{quote}

That all the $\chi_1(p)$ are roots of unity of degree $f_1$ is
clear, for it is always the case (cf. \S.\,177) that
$p^{f_1}\equiv\pm1\pmod{m}$, and therefore
\[ [\chi_1(p)]^{f_1} = \chi_1(\pm1) = 1. \]
But if a primitive $f_1$th root of unity occurs among the
$\chi_1(p)$, then all the other $f_1$th roots of unity occur among
them as well, and equally often, namely as often as the root of
unity 1.

This follows immediately from the group nature of the $\chi_1$. It
therefore remains only to show that \gesperrt{a primitive} $f_1$th root of
unity occurs among the $\chi_1(p)$.

But this possibility follows very simply from the determination of
$\chi_1$. Suppose first that, for \gesperrt{odd $m$, $\Theta$ in (6) is a
primitive root of unity of degree} $\tfrac12\mu$, and that $\gamma$
is the index of $p$; then a power of $\Theta^{\gamma}$,
$\Theta^{\gamma\lambda}$, is $=1$ then and only then, when
\[ 2\gamma\lambda \equiv 0 \pmod{\mu}. \]
Now $\mu=ef$, and $e$ is the greatest common divisor of $\gamma$
and $\mu$; and therefore this condition becomes
$2\lambda\equiv0\pmod{f}$. But by \S.\,177, 3. $f_1=f$
or $=\tfrac12 f$, according as $f$ is odd or even, and
consequently $\lambda\equiv0\pmod{f_1}$ must hold. Thus
$\Theta^{\gamma}$ is a primitive $f_1$th root of unity.

But if $m$ is \gesperrt{even}, then in (7) we likewise take for $\Theta$
a primitive $\tfrac12\mu$th root of unity and put
\[ p \equiv (-1)^{\alpha}5^{\beta} \pmod{m}, \]
so that $f$ is the least positive number for which the congruences
\[ \alpha f \equiv 0 \pmod{2}, \quad \beta f \equiv 0 \pmod{\tfrac12\mu} \]
are satisfied. Now $\Theta^{\beta\lambda}$ is $=1$ only when
$\beta\lambda\equiv0\pmod{\tfrac12\mu}$. Unless $\lambda=1$
(which forces $\beta=0$), the least possible value of $\lambda$ is
even and hence equals $f$.
